The encoding described in this chapter defines the architectural instruction
encoding for VIV-32 v\version{}.

\section{Instruction Word}

\begin{itemize}
\item VIV-32 instructions are fixed-width 32-bit words.
\item Instruction fetch addresses must be 4-byte aligned.
\item Instruction words are stored in memory using the architecture's big-endian
      byte order.
\item Instruction encoding diagrams number bits from bit 31, the most
      significant bit, to bit 0, the least significant bit.
\item Unless otherwise stated, immediate fields are sign-extended to 32 bits
      before use.
\item Logical immediate instructions zero-extend their immediate operands.
\end{itemize}

For example, the instruction word \code{0x12345678} stored at address \code{A}
is represented in memory as follows:

\begin{center}
\begin{tabular}{ll}
\hline
Address & Byte value \\
\hline
\code{A + 0} & \code{0x12} \\
\code{A + 1} & \code{0x34} \\
\code{A + 2} & \code{0x56} \\
\code{A + 3} & \code{0x78} \\
\hline
\end{tabular}
\end{center}

\section{Major Opcode Field}

\begin{itemize}
\item Bits 31--26 contain the major opcode.
\item The major opcode is 6 bits wide.
\item A 6-bit major opcode provides 64 primary opcode classes.
\item Individual opcode classes may define additional function fields.
\item Opcode space is reserved for future extensions.
\end{itemize}

\begin{center}
\begin{tabular}{lll}
\hline
Bits & Field & Description \\
\hline
31--26 & opcode & Major opcode \\
25--0  & payload & Format-specific instruction payload \\
\hline
\end{tabular}
\end{center}

\section{Register Fields}

\begin{itemize}
  \item VIV-32 v\version{} has 16 general-purpose registers, and 8
        floating-point registers.
  \item Register fields are 4 bits wide.
  \item A register field value of \code{0x0} selects register \code{\$0} or
        register \code{\#0}, depending on the instruction context.
  \item A register field value of \code{0xF} selects GPR register \code{\$15},
        but is invalid for the FPR as it only goes to \code{\#7}.
  \item Register field names such as \code{rd}, \code{ra}, \code{rb},
        \code{rd0}, \code{rd1}, \code{rs}, \code{base}, and \code{target} are
        documentation names for encoded register fields.
  \item Register fields are placed on nibble boundaries. Register fields are
        placed in the least-significant nibbles of the instruction word to aid
        in reading.
\end{itemize}

\begin{center}
\begin{tabular}{llc}
\hline
Field value & GPR & FPR\\
\hline
\code{0x0} & \code{\$0} & \code{\#0} \\
\code{0x1} & \code{\$1} & \code{\#1} \\
\code{0x2} & \code{\$2} & \code{\#2} \\
\code{0x3} & \code{\$3} & \code{\#3} \\
\code{0x4} & \code{\$4} & \code{\#4} \\
\code{0x5} & \code{\$5} & \code{\#5} \\
\code{0x6} & \code{\$6} & \code{\#6} \\
\code{0x7} & \code{\$7} & \code{\#7} \\
\code{0x8} & \code{\$8} & - \\
\code{0x9} & \code{\$9} & - \\
\code{0xA} & \code{\$10} & - \\
\code{0xB} & \code{\$11} & - \\
\code{0xC} & \code{\$12} & - \\
\code{0xD} & \code{\$13} & - \\
\code{0xE} & \code{\$14} & - \\
\code{0xF} & \code{\$15} & - \\
\hline
\end{tabular}
\end{center}

\section{Instruction Format Families}

VIV-32 v\version{} defines a small number of instruction format families. The
exact opcode assignments are defined separately in the opcode map.

\subsection{R-Type: Register-Register Operations}

R-type instructions operate primarily on registers. Register fields are placed
in the least-significant nibbles of the instruction word so that register
operands remain readable in hexadecimal instruction dumps.

\begin{center}
\begin{tabular}{lllllll}
\hline
Bits & 31--26 & 25--24 & 23--12 & 11--8 & 7--4 & 3--0 \\
\hline
Field & \code{opcode} & \code{file} & \code{func/reserved} & \code{rd} &
\code{ra} & \code{rb} \\
\hline
\end{tabular}
\end{center}

\begin{itemize}
  \item \code{opcode} is the 6-bit major opcode.
  \item \code{file} is the target register file---GPR or FPR. Additional file
        values have been reserved for future additional register files, such
        as SIMD.
  \item \code{func/reserved} selects the specific register-register operation or
        is reserved.
  \item \code{rd} is the destination register field.
  \item \code{ra} is the first source register field.
  \item \code{rb} is the second source register field.
  \item Arithmetic and logical R-type instructions update condition flags by
        default.
\end{itemize}

\noindent
Examples of R-type operations include:

\begin{itemize}
  \item \code{add \$rd, \$ra, \$rb}
  \item \code{subd \#rd, \#ra, \#rb}
  \item \code{and \$rd, \$ra, \$rb}
  \item \code{or \$rd, \$ra, \$rb}
  \item \code{xord \#rd, \#ra, \#rb}
  \item \code{itod \#rd, \$ra ; implicit \$rb=\$0} 
\end{itemize}

\subsection{R2-Type: Register-Register Operations with Two Destinations}

R2-type instructions are used for operations that naturally produce two register
results, such as full-width multiplication or division with remainder. The four
register fields occupy the least-significant four nibbles of the instruction
word.

\begin{center}
\begin{tabular}{llllllll}
\hline
Bits & 31--26 & 25--24 & 23--16 & 15--12 & 11--8 & 7--4 & 3--0 \\
\hline
Field & \code{opcode} & \code{file} & \code{func/reserved} & \code{rd0} &
\code{rd1} & \code{ra} & \code{rb} \\
\hline
\end{tabular}
\end{center}

\begin{itemize}
  \item \code{opcode} is the 6-bit major opcode.
  \item \code{file} is the target register file---GPR or FPR.
  \item \code{func/reserved} selects the specific two-result operation or is
        reserved.
  \item \code{rd0} is the first destination register field.
  \item \code{rd1} is the second destination register field. Note that for FPR
        instructions, this should be \code{0}.
  \item \code{ra} is the first source register field.
  \item \code{rb} is the second source register field.
\end{itemize}

\noindent
Examples of R2-type operations include:

\begin{itemize}
  \item \code{mul \$rdlo, \$rdhi, \$ra, \$rb}
  \item \code{muld \#rdlo, \#0, \#ra, \#rb}
  \item \code{divd \#rdq, \#0, \#ra, \#rb}
  \item \code{divu \$rdq, \$rdr, \$ra, \$rb}
\end{itemize}

\noindent
For GPR multiplication instructions, \code{\$rd0} receives the low 32 bits of
the result and \code{\$rd1} receives the high 32 bits of the result.
For GPR division instructions, \code{\$rd0} receives the quotient and
\code{\$rd1} receives the remainder.

\subsection{I-Type: Register-Immediate Operations}

I-type instructions operate on a GPR register and a 16-bit immediate value.
Register fields are placed in the least-significant nibbles of the instruction
word so that the register operands remain readable in hexadecimal dumps.

\begin{center}
\begin{tabular}{llllll}
\hline
Bits & 31--26 & 25--24 & 23--8 & 7--4 & 3--0 \\
\hline
Field & \code{opcode} & \code{mode/reserved} & \code{imm16} & \code{rd} &
\code{ra} \\
\hline
\end{tabular}
\end{center}

\begin{itemize}
  \item \code{opcode} is the 6-bit major opcode.
  \item \code{mode/reserved} selects an instruction variant or is reserved.
  \item \code{imm16} is a 16-bit immediate field.
  \item \code{rd} is the destination register field.
  \item \code{ra} is the source register field.
  \item Unless otherwise stated, \code{imm16} is sign-extended to 32 bits
        before use.
  \item Logical immediate instructions zero-extend \code{imm16} to 32 bits
        before use.
  \item Arithmetic and logical I-type instructions update condition flags by
        default.
\end{itemize}

Examples of I-type operations include:

\begin{itemize}
  \item \code{addi \$rd, \$ra, imm}
  \item \code{subi \$rd, \$ra, imm}
  \item \code{andi \$rd, \$ra, imm}
  \item \code{ori \$rd, \$ra, imm}
  \item \code{xori \$rd, \$ra, imm}
  \item \code{cmpi \$ra, imm}
\end{itemize}

\subsection{BI-Type: Register-Immediate Bit Operations}

BI-type instructions operate on a register and a 14-bit immediate value.
Register fields are placed in the least-significant nibbles of the
instruction word so that the register operands remain readable in
hexadecimal dumps.

\begin{center}
\begin{tabular}{lllllll}
\hline
Bits & 31--26 & 25--24 & 23--22 & 21--8 & 7--4 & 3--0 \\
\hline
Field & \code{opcode} & \code{mode/reserved} & \code{file} & \code{imm14} &
\code{rd} & \code{ra} \\
\hline
\end{tabular}
\end{center}

\begin{itemize}
  \item \code{opcode} is the 6-bit major opcode.
  \item \code{mode/reserved} selects an instruction variant or is reserved.
  \item \code{file} is the target register file---GPR or FPR.
  \item \code{imm14} is a 14-bit immediate field.
  \item \code{rd} is the destination register field.
  \item \code{ra} is the source register field.
  \item Arithmetic and logical BI-type instructions update condition flags by
        default.
\end{itemize}

Examples of BI-type operations include:

\begin{itemize}
  \item \code{btst \$rd, \$ra, imm}
  \item \code{bsetd \#rd, \#ra, imm}
\end{itemize}

\subsection{U-Type: Upper Immediate and Constant Construction}

U-type instructions construct constants and addresses in general-purpose
registers. U-type instructions use a 16-bit immediate field and place the
destination register in the least-significant nibble of the instruction word.

\begin{center}
\begin{tabular}{lllll}
\hline
Bits & 31--26 & 25--20 & 19--4 & 3--0 \\
\hline
Field & \code{opcode} & \code{mode/reserved} & \code{imm16} & \code{rd} \\
\hline
\end{tabular}
\end{center}

\begin{itemize}
  \item \code{opcode} is the 6-bit major opcode.
  \item \code{mode/reserved} selects the constant-construction operation or is
        reserved.
  \item \code{imm16} is a 16-bit immediate field.
  \item \code{rd} is the destination register field.
  \item U-type instructions are intended to support full 32-bit constant and
        address construction.
  \item U-type immediates are interpreted according to the specific instruction.
\end{itemize}

\noindent
Examples of U-type operations include:

\begin{itemize}
  \item \code{lui \$rd, imm16}
  \item \code{lli \$rd, imm16}
  \item \code{lhi \$rd, imm16}
\end{itemize}

\noindent
A typical full-width constant may be constructed using a lower-immediate
instruction followed by a high-immediate instruction:

\begin{verbatim}
    lli  $1, 0x5678
    lhi  $1, 0x1234
\end{verbatim}

\noindent
This sequence of instructions is also represented by the pseudo-instruction
\code{li}.

\subsection{M-Type: Load and Store Operations}

M-type instructions access memory using base-plus-offset addressing. The
register fields are placed in the least-significant nibbles of the instruction
word so that the load/store register and base register remain readable in
hexadecimal dumps.

\begin{center}
\begin{tabular}{lllllll}
\hline
Bits & 31--26 & 25--24 & 23 & 22--8 & 7--4 & 3--0 \\
\hline
Field & \code{opcode} & \code{size} & \code{mode} & \code{offset15} &
\code{rd/rs} & \code{base} \\
\hline
\end{tabular}
\end{center}

\begin{itemize}
  \item \code{opcode} is the 6-bit major opcode. Load and store instructions use
        separate major opcodes.
  \item \code{size} encodes the width of the operation.
  \item \code{mode} flags extra load properties, determining if the operation is
        signed, unsigned or float, relative to the \code{size}.
  \item \code{offset15} is a signed 15-bit byte offset.
  \item \code{rd/rs} is the destination register field for loads.
  \item \code{rd/rs} is the source register field for stores.
  \item \code{base} is the base address register field.
  \item The effective address is computed as
        \ensuremath{base + \sext{offset15}}.
\end{itemize}

\subsubsection{Load Interpretation}

For load instructions, \code{mode} selects whether sub-word loads are
zero-extended for an unsigned load, sign-extended for a signed load; or if a
word load is a binary32; or if a double word load's target is the upper or
lower word;

\begin{center}
\begin{tabular}{lllll}
\hline
Mnemonic & \code{size} & \code{mode} & Meaning & File \\
\hline
\code{lbu} & \code{00} & \code{0} & load unsigned byte     & GPR \\
\code{lb}  & \code{00} & \code{1} & load signed byte       & GPR \\
\code{lhu} & \code{01} & \code{0} & load unsigned halfword & GPR \\
\code{lh}  & \code{01} & \code{1} & load signed halfword   & GPR \\
\code{lw}  & \code{10} & \code{0} & load word              & GPR \\
\code{ld}  & \code{10} & \code{1} & load binary32          & FPR \\
\code{ldl} & \code{11} & \code{0} & load binary64 lower    & FPR \\
\code{ldu} & \code{11} & \code{1} & load binary64 upper    & FPR \\  
\hline
\end{tabular}
\end{center}

\subsubsection{Store Interpretation}

For store instruction, \code{mode} only applies to word-sized operations,
determining if the target is a float, and which register.

\begin{center}
\begin{tabular}{lllll}
\hline
Mnemonic & \code{size} & \code{mode} & Meaning & File \\
\hline
\code{sb}  & \code{00} & \code{0} & store byte           & GPR \\
\code{sh}  & \code{01} & \code{0} & store halfword       & GPR \\
\code{sw}  & \code{10} & \code{0} & store word           & GPR \\
\code{sd}  & \code{10} & \code{1} & store binary32       & FPR \\
\code{sdl} & \code{11} & \code{0} & store lower binary64 & FPR \\
\code{sdu} & \code{11} & \code{1} & store upper binary64 & FPR \\
\hline
\end{tabular}
\end{center}

\subsubsection{Addressing Examples}

The \code{rd/rs} registers belong to the file targeted by the instruction. The
base registers are always from the GPR. Examples of M-type load operations
include:

\begin{itemize}
  \item \code{lb \$rd, [\$base, offset]}
  \item \code{lbu \$rd, [\$base, offset]}
  \item \code{lh \$rd, [\$base, offset]}
  \item \code{ldl \#rd, [\$base, offset]}
  \item \code{ld \#rd, [\$base, offset]}
\end{itemize}

\noindent
Examples of M-type store operations include:

\begin{itemize}
  \item \code{sb \$rs, [\$base, offset]}
  \item \code{sh \$rs, [\$base, offset]}
  \item \code{sdu \#rs, [\$base, offset]}
\end{itemize}

\subsubsection{Floating-point Expansion}

The floating-point instructions operate only on the floating-point registers.
As the registers are 64-bit, binary32 loads expand their IEEE-754 binary32
value to binary64 representation, and binary32 stores convert the register's
binary64 value to binary32.

Floating-point loads update the floating-point condition flags according to
the resulting value. In particular, \code{N} and \code{Z} describe the
numeric result, while \code{I}, \code{U}, \code{X}, and \code{E} indicate
the corresponding floating-point result or conversion conditions. Integer-only
flags \code{C} and \code{V} are cleared.

The binary64 half-load and half-store operations transfer raw 32-bit portions
of a floating-point register without converting or otherwise interpreting the
transferred portion. A half-load modifies only the targeted 32 bits of the
destination floating-point register; the non-targeted portion is preserved.

After each half-load, the condition flags are updated from the complete
64-bit contents of the destination floating-point register as though those
contents were an IEEE-754 binary64 value. Consequently, when constructing a
binary64 value using two half-loads, the flag state produced by the first
half-load is generally not meaningful. The flags should be considered to
describe the loaded binary64 value only after both halves have been loaded.

The following stores and then loads a binary64 value without loss of
precision:

\begin{verbatim}
sdu     #1, [$1, 0]
sdl     #1, [$1, 4]
ldu     #2, [$1, 0]
ldl     #2, [$1, 4]
\end{verbatim}

Unlike the U-type instructions, these operations do \textbf{not} modify the
non-targeted portion of the floating-point register. The order of the two
half operations therefore does not affect the final register value, although
it may affect the transient condition flags produced between the two loads.

These operation sequences also form the pseudo-instructions \code{sdd} and
\code{ldd}. Note that the ordering here is deliberately big-endian to match
the rest of the architecture.

\subsection{BF-Type: Flag-Based Conditional Branches}

BF-type instructions branch according to the current condition flags in the
status register. They test a single existing status flag and do not perform
arithmetic, comparison, or interpretation of the flag state.

BF-type instructions do not contain register fields, so the condition field
and branch offset occupy the instruction payload.

\begin{center}
\begin{tabular}{llll}
\hline
Bits & 31--26 & 25--22 & 21--0 \\
\hline
Field & \code{opcode} & \code{cond} & \code{offset22} \\
\hline
\end{tabular}
\end{center}

\begin{itemize}
\item \code{opcode} is the 6-bit major opcode.
\item \code{cond} is the 4-bit flag-condition field.
\item \code{offset22} is a signed 22-bit PC-relative branch offset.
\item Branch offsets are measured in 32-bit instruction words.
\item The encoded offset is shifted left by two bits to form a byte offset.
\item The branch target is relative to the address of the instruction following
the branch.
\item BF-type instructions test existing condition flags; they do not perform a
new comparison or otherwise interpret the flags.
\item BF-type instructions do not update condition flags.
\end{itemize}

The \code{cond} field identifies one of the eight architectural condition
flags, together with whether the branch tests for that flag being set or
clear. The available flags are \code{N}, \code{Z}, \code{C}, \code{V},
\code{E}, \code{I}, \code{U}, and \code{X}.

\noindent
Examples of BF-type operations include:

\begin{center}
\begin{tabular}{ll}
\hline
Instruction & Meaning \\
\hline
\code{bf.ns label} & Branch if \code{N} is set. \\
\code{bf.nc label} & Branch if \code{N} is clear. \\
\code{bf.zs label} & Branch if \code{Z} is set. \\
\code{bf.zc label} & Branch if \code{Z} is clear. \\
\code{bf.cs label} & Branch if \code{C} is set. \\
\code{bf.cc label} & Branch if \code{C} is clear. \\
\code{bf.vs label} & Branch if \code{V} is set. \\
\code{bf.vc label} & Branch if \code{V} is clear. \\
\code{bf.es label} & Branch if \code{E} is set. \\
\code{bf.ec label} & Branch if \code{E} is clear. \\
\code{bf.is label} & Branch if \code{I} is set. \\
\code{bf.ic label} & Branch if \code{I} is clear. \\
\code{bf.us label} & Branch if \code{U} is set. \\
\code{bf.uc label} & Branch if \code{U} is clear. \\
\code{bf.xs label} & Branch if \code{X} is set. \\
\code{bf.xc label} & Branch if \code{X} is clear. \\
\hline
\end{tabular}
\end{center}

BF-type branches have no signed, unsigned, integer, or floating-point
comparison semantics. Comparisons requiring interpretation of register values
are performed using the register-comparison branch instructions.

\subsection{BC-Type: Register-Comparison Branches}

BC-type instructions compare two registers and branch according to the
specified comparison. The register fields are placed in the least-significant
nibbles of the instruction word so that the compared registers remain readable
in hexadecimal instruction dumps.

\begin{center}
\begin{tabular}{llllll}
\hline
Bits & 31--26 & 25--22 & 21--8 & 7--4 & 3--0 \\
\hline
Field & \code{opcode} & \code{cond} & \code{offset14} & \code{ra} & \code{rb} \\
\hline
\end{tabular}
\end{center}

\begin{itemize}
  \item \code{opcode} is the 6-bit major opcode. This will determine if the
        operation is GPR or FPR. Note that operation codes in the GPR
        instruction may not map to the same operation in the FPR instruction.
  \item \code{cond} is the 4-bit comparison condition-code field.
  \item \code{offset14} is a signed 14-bit PC-relative branch offset.
  \item \code{ra} is the first comparison register field.
  \item \code{rb} is the second comparison register field.
  \item Branch offsets are measured in 32-bit instruction words.
  \item The encoded offset is shifted left by two bits to form a byte offset.
  \item The branch target is relative to the address of the instruction
        following the branch.
  \item BC-type instructions compare \code{ra} and \code{rb} directly.
  \item BC-type instructions do not depend on existing condition flags.
  \item BC-type instructions do not update condition flags.
\end{itemize}

\noindent
The low byte of a BC-type instruction contains the compared register fields.
This makes the two comparison operands directly visible in the
least-significant two hexadecimal digits of the instruction word.

\noindent
Examples of BC-type operations include:

\begin{center}
\begin{tabular}{ll}
\hline
Instruction & Meaning \\
\hline
\code{b.eq  \$ra, \$rb, label}  & Branch to \code{label} if \code{ra} equals
\code{rb}. \\
\code{b.ne  \$ra, \$rb, label}  & Branch to \code{label} if \code{ra} does not equal
\code{rb}. \\
\code{bd.lt  \#ra, \#rb, label}  & Branch to \code{label} if \code{ra} is less than
\code{rb}. \\
\code{b.ge  \$ra, \$rb, label}  &
\begin{tabular}[t]{@{}l@{}}
    Branch to \code{label} if \code{ra} is greater than or equal to \code{rb}, treating\\
    both as signed values.
\end{tabular} \\
\code{b.ltu \$ra, \$rb, label} &
\begin{tabular}[t]{@{}l@{}}
    Branch to \code{label} if \code{ra} is less than \code{rb}, treating both as
    unsigned \\ values.
\end{tabular} \\
\code{bd.ge \#ra, \#rb, label} & Branch to \code{label} if \code{ra} is greater than
or equal to \code{rb}. \\
\hline
\end{tabular}
\end{center}

\subsection{J-Type: PC-Relative Jump and Call}

J-type instructions perform unconditional PC-relative control transfer. J-type
instructions do not contain register fields, so the full instruction payload is
available for the PC-relative offset.

\begin{center}
\begin{tabular}{lll}
\hline
Bits & 31--26 & 25--0 \\
\hline
Field & \code{opcode} & \code{offset26} \\
\hline
\end{tabular}
\end{center}

\begin{itemize}
  \item \code{opcode} is the 6-bit major opcode.
  \item \code{offset26} is a signed 26-bit PC-relative offset.
  \item Jump offsets are measured in 32-bit instruction words.
  \item The encoded offset is shifted left by two bits to form a byte offset.
  \item The jump target is relative to the address of the instruction following
        the jump or call instruction.
  \item J-type instructions do not update condition flags.
\end{itemize}

\noindent
Examples of J-type operations include:

\begin{itemize}
  \item \code{jmp label}
  \item \code{call label}
\end{itemize}

\noindent
The \code{jmp} instruction changes the program counter to the computed target
address. The \code{call} instruction also writes the return address to the
architectural link register, \code{\$15}.

\subsection{JR-Type: Register-Indirect Jump and Call}

JR-type instructions perform register-indirect control transfer. The register
fields are placed in the least-significant nibbles of the instruction word so
that the link register and target register remain readable in hexadecimal dumps.

\begin{center}
\begin{tabular}{lllll}
\hline
Bits & 31--26 & 25--8 & 7--4 & 3--0 \\
\hline
Field & \code{opcode} & \code{func/reserved} & \code{rd/link} & \code{target}\\
\hline
\end{tabular}
\end{center}

\begin{itemize}
  \item \code{opcode} is the 6-bit major opcode.
  \item \code{func/reserved} selects the register-indirect control-transfer
        operation or is reserved.
  \item \code{rd/link} is the destination link register field for
        register-indirect calls.
  \item \code{target} is the target address register field.
  \item For plain register jumps, \code{rd/link} must be encoded as \code{\$0}.
  \item For register-indirect calls, \code{rd/link} receives the return address.
  \item The standard ABI link register is \code{\$15}.
  \item Register-indirect targets must be 4-byte aligned.
  \item JR-type instructions do not update condition flags.
\end{itemize}

\noindent
Examples of JR-type operations include:

\begin{itemize}
  \item \code{jr \$target}
  \item \code{jalr \$rd, \$target}
  \item \code{ret}
\end{itemize}

\noindent
In VIV-32 v\version{}, \code{ret} is defined as a pseudo-instruction for
\code{jr \$15}.

\subsection{X-Type: System and Control Operations}

X-type instructions access system facilities, control registers, traps,
interrupt state, and processor execution state. The register fields are placed
in the least-significant nibbles of the instruction word so that source and
destination registers remain readable in hexadecimal dumps.

\begin{center}
\begin{tabular}{lllll}
\hline
Bits & 31--26 & 25--8 & 7--4 & 3--0 \\
\hline
Field & \code{opcode} & \code{sysfunc} & \code{rd} & \code{rs} \\
\hline
\end{tabular}
\end{center}

\begin{itemize}
  \item \code{opcode} is the 6-bit major opcode.
  \item \code{sysfunc} is an 18-bit field divided into \code{sysop} and
        \code{payload}.
  \item \code{sysop} selects the system or control operation.
  \item \code{payload} supplies operation-specific data, such as a trap
        immediate or control-register number.  
  \item \code{rd} is the destination register field when required.
  \item \code{rs} is the source register field when required.
  \item Instructions that do not require \code{\$rd} or \code{\$rs} must encode
        unused register fields as \code{\$0}.
  \item X-type instructions do not update condition flags unless explicitly
        specified by the instruction set chapter.
\end{itemize}

\noindent
Examples of X-type operations include:

\begin{itemize}
  \item \code{nop}
  \item \code{halt}
  \item \code{trap imm}
  \item \code{syscall}
  \item \code{iret}
  \item \code{ei}
  \item \code{di}
  \item \code{rdpc \$rd}
  \item \code{mrs \$rd, \%cr}
  \item \code{msr \%cr, \$rs}
\end{itemize}

\noindent
For control-register access instructions, the selected control register is
encoded in the \code{creg4} field within the X-type \code{payload}. The
\code{\$rd} and \code{\$rs} fields are used only for the general-purpose
register operand.

The detailed interpretation of \code{sysfunc}, including \code{sysop},
\code{payload}, \code{trapimm12}, and \code{creg4}, is defined in the system
and control encoding section.

\section{Major Opcode Map}

The 6-bit \code{opcode} field identifies the primary instruction class. Most
instruction classes use additional fields, such as \code{func}, \code{mode},
\code{cond}, or \code{sysfunc}, to select the specific instruction within that
class.

\begin{center}
\begin{tabular}{llll}
\hline
Opcode & Format & Class & Instructions \\
\hline
\code{0x00} & X & System/control &
\begin{tabular}[t]{@{}l@{}}
\code{nop}, \code{halt}, \code{trap}, \code{syscall}, \code{iret}, \code{ei},\\
\code{di}, \code{rdpc}, \code{mrs}, \code{msr}
\end{tabular} \\

\code{0x01} & R  & Register ALU &
\begin{tabular}[t]{@{}l@{}}
\code{add}, \code{sub}, \code{and}, \code{or}, \code{xor}, \code{not},
\code{neg}, \code{cmp}, \\
\code{addd}, \code{subd}, \code{andd}, \code{ord}, \code{xord}, \code{notd}, \\
\code{negd}, \code{cmpd}
\end{tabular} \\

\code{0x02} & I  & Immediate arithmetic/compare & \code{addi}, \code{subi},
\code{cmpi} \\

\code{0x03} & I  & Immediate logical & \code{andi}, \code{ori}, \code{xori} \\

\code{0x04} & R  & Register shifts & \code{shl}, \code{shr}, \code{sar},
\code{shld}, \code{shrd}, \code{sard} \\

\code{0x05} & I  & Immediate shifts & \code{shli}, \code{shri}, \code{sari} \\

\code{0x06} & BI  & Bit immediate &
\begin{tabular}[t]{@{}l@{}}
\code{shldi}, \code{shrdi}, \code{sardi}, \code{btst}, \code{btstd},\\
\code{bset}, \code{bsetd}, \code{bclr}, \code{bclrd}, \code{btgl}, \\
\code{btgld}
\end{tabular} \\

\code{0x07} & R2 & Multiply/divide & \code{mul}, \code{mulu}, \code{div},
\code{divu}, \code{muld}, \code{divd} \\

\code{0x08} & U  & Constant construction & \code{lui}, \code{lli}, \code{lhi}\\

\code{0x09} & M  & Load & \code{lb}, \code{lbu}, \code{lh}, \code{lhu},
\code{lw}, \code{ld}, \code{ldl}, \code{ldu} \\

\code{0x0A} & M  & Store & \code{sb}, \code{sh}, \code{sw}, \code{sd},
\code{sdl}, \code{sdu} \\

\code{0x0B} & BF & Flag branch & \code{bf.*} \\

\code{0x0C} & BC & GPR branch & \code{b.*} \\

\code{0x0D} & J  & PC-relative jump & \code{jmp} \\

\code{0x0E} & J  & PC-relative call & \code{call} \\

\code{0x0F} & JR & Register jump/call & \code{jr}, \code{jalr} \\

\code{0x10} & R & Inter-file transfer & \code{disf}, \code{itod}, \code{utod}, \code{ftod},
\code{dtoi}, \code{dtou}, \code{dtof} \\

\code{0x11} & BC & FPR branch & \code{bd.*} \\

\code{0x12--0x3F} & -- & Reserved & Reserved for future extensions \\
\hline
\end{tabular}
\end{center}

\noindent
Opcode values not listed as assigned are reserved. Reserved opcode values must
not be emitted by conforming assemblers. Executing an instruction with a
reserved opcode raises an \code{IllegalInstructionException}.

\section{Class-Specific Encoding Values}

This section defines the class-specific function and mode fields used by the
instruction classes listed in the major opcode map.

\subsection{Register ALU Function Codes}

Register ALU instructions use opcode \code{0x01} and the R-type instruction
format.

\begin{center}
\begin{tabular}{llll}
\hline
\code{file} & \code{func} value & Mnemonic & Meaning \\
\hline
\code{0x0} & \code{0x000} & \code{add} & Add registers, GPR \\
\code{0x1} & \code{0x000} & \code{addd} & Add registers, FPR \\
\code{0x0} & \code{0x001} & \code{sub} & Sub registers, GPR \\
\code{0x1} & \code{0x001} & \code{subd}  & Sub registers, FPR \\
\code{0x0} & \code{0x002} & \code{and} & Bitwise AND, GPR \\
\code{0x1} & \code{0x002} & \code{andd} & Bitwise AND, FPR \\
\code{0x0} & \code{0x003} & \code{or} & Bitwise OR, GPR \\
\code{0x1} & \code{0x003} & \code{ord} & Bitwise OR, FPR \\
\code{0x0} & \code{0x004} & \code{xor} & Bitwise XOR, GPR \\
\code{0x1} & \code{0x004} & \code{xord} & Bitwise XOR, FPR \\
\code{0x0} & \code{0x005} & \code{not} & Bitwise NOT, GPR \\
\code{0x1} & \code{0x005} & \code{notd}  & Bitwise NOT, FPR \\
\code{0x0} & \code{0x006} & \code{neg} & Negate register, GPR \\
\code{0x1} & \code{0x006} & \code{negd} & Negate register, FPR \\
\code{0x0} & \code{0x007} & \code{cmp} & Compare registers, GPR \\
\code{0x1} & \code{0x007} & \code{cmpd} & Compare registers, FPR \\
\code{0x0--0x3} & \code{0x008--0xFFF} & reserved & Reserved \\
\hline
\end{tabular}
\end{center}

\noindent
For \code{not}, \code{notd}, \code{neg}, \code{negd}, \code{cmp}, \code{cmpd}
the \code{rb} field is reserved and must be encoded as \code{0} by conforming
assemblers.

\subsection{Immediate Arithmetic and Compare Mode Codes}

Immediate arithmetic and compare instructions use opcode \code{0x02} and the
I-type instruction format.

\begin{center}
\begin{tabular}{lll}
\hline
\code{mode} value & Mnemonic & Meaning \\
\hline
\code{0x0} & \code{addi} & Add sign-extended immediate \\
\code{0x1} & \code{subi} & Subtract sign-extended immediate \\
\code{0x2} & \code{cmpi} & Compare with sign-extended immediate \\
\code{0x3} & reserved & Reserved \\
\hline
\end{tabular}
\end{center}

\noindent
For \code{cmpi}, the \code{rd} field is reserved and must be encoded as \code{r0} by conforming assemblers. The instruction updates condition flags but does not write a general-purpose register.

\subsection{Immediate Logical Mode Codes}

Immediate logical instructions use opcode \code{0x03} and the I-type instruction format.

\begin{center}
\begin{tabular}{lll}
\hline
\code{mode} value & Mnemonic & Meaning \\
\hline
\code{0x0} & \code{andi} & Bitwise AND with zero-extended immediate \\
\code{0x1} & \code{ori}  & Bitwise OR with zero-extended immediate \\
\code{0x2} & \code{xori} & Bitwise XOR with zero-extended immediate \\
\code{0x3} & reserved & Reserved \\
\hline
\end{tabular}
\end{center}

\subsection{Register Shift Function Codes}

Register shift instructions use opcode \code{0x04} and the R-type instruction
format.

\begin{center}
\begin{tabular}{llll}
\hline
\code{file} & \code{func} value & Mnemonic & Meaning \\
\hline
\code{0x0} & \code{0x000} & \code{shl} & Logical shift left by register count, GPR \\
\code{0x1} & \code{0x000} & \code{shld} & Logical shift left by register count, FPR \\
\code{0x0} & \code{0x001} & \code{shr} & Logical shift right by register count, GPR \\
\code{0x1} & \code{0x001} & \code{shrd} & Logical shift right by register count, FPR \\
\code{0x0} & \code{0x002} & \code{sar} & Arithmetic shift right by register count, GPR \\
\code{0x1} & \code{0x002} & \code{sard} & Arithmetic shift right by register count, FPR \\
\code{0x0--0x3} & \code{0x003--0xFFF} & reserved & Reserved \\
\hline
\end{tabular}
\end{center}

\subsection{Immediate Shift Mode Codes}

Immediate shift instructions use opcode \code{0x05} and the I-type instruction format.

\begin{center}
\begin{tabular}{lll}
\hline
\code{mode} value & Mnemonic & Meaning \\
\hline
\code{0x0} & \code{shli} & Logical shift left by immediate count \\
\code{0x1} & \code{shri} & Logical shift right by immediate count \\
\code{0x2} & \code{sari} & Arithmetic shift right by immediate count \\
\code{0x3} & reserved & Reserved \\
\hline
\end{tabular}
\end{center}

\subsection{Bit Immediate Mode Codes}

Bit immediate instructions use opcode \code{0x06} and the BI-type instruction format.

\begin{center}
\begin{tabular}{llll}
\hline
\code{mode} &\code{file} & Mnemonic & Meaning \\
\hline
\code{0x0} & \code{0x0} & \code{btst} & Test immediate bit index, GPR \\
\code{0x0} & \code{0x1} & \code{btstd} & Test immediate bit index, FPR \\
\code{0x1} & \code{0x0} & \code{bset} & Set immediate bit index, GPR \\
\code{0x1} & \code{0x1} & \code{bsetd} & Set immediate bit index, FPR \\
\code{0x2} & \code{0x0} & \code{bclr} & Clear immediate bit index, GPR \\
\code{0x2} & \code{0x1} & \code{bclrd} & Clear immediate bit index, FPR \\
\code{0x3} & \code{0x0} & \code{btgl} & Toggle immediate bit index, GPR \\
\code{0x3} & \code{0x1} & \code{btgld} & Toggle immediate bit index, FPR \\
\code{0x4} & \code{0x1} & \code{shldi} & Logical shift left, FPR \\
\code{0x5} & \code{0x1} & \code{shrdi} & Logical shift right, FPR \\
\code{0x6} & \code{0x1} & \code{sardi} & Arithmetic shift right, FPR \\
\hline
\end{tabular}
\end{center}

\noindent
The value in \code{imm14} is the bit index or shift count, according to the
instruction.

\noindent
For \code{btst}, and \code{btstd} the \code{rd} field is reserved and must be
encoded as \code{0} by conforming assemblers. The instruction updates condition
flags but does not write to a register.

\subsection{Multiply and Divide Function Codes}

Multiply and divide instructions use opcode \code{0x07} and the R2-type
instruction format.

\begin{center}
\begin{tabular}{llll}
\hline
\code{file} & \code{func} value & Mnemonic & Meaning \\
\hline
\code{0} & \code{0x00} & \code{mul}  & Signed full-width multiply, GPR \\
\code{1} & \code{0x00} & \code{muld} & Floating-point multiply, FPR \\
\code{0} & \code{0x01} & \code{mulu}  & Unsigned full-width multiply, GPR \\
\code{1} & \code{0x01} & - & Reserved \\
\code{0} & \code{0x02} & \code{div}  & Signed divide with remainder, GPR \\
\code{1} & \code{0x02} & \code{divd} & Floating-point divide, FPR \\
\code{0} & \code{0x03} & \code{divu}  & Unsigned divide with remainder, GPR \\
\code{1} & \code{0x03} & - & Reserved \\
\code{-} & \code{0x04--0xFF} & reserved & Reserved \\
\hline
\end{tabular}
\end{center}

\noindent
For GPR multiply instructions, \code{rd0} receives the low 32 bits of the
product and \code{rd1} receives the high 32 bits of the product. FPR multiply
results are stored in a single register, with \code{rd1} being \code{0}.

\noindent
For GPR divide instructions, \code{rd0} receives the quotient and \code{rd1}
receives the remainder. FPR divide results are stored in a single register,
with \code{rd1} being \code{0}.

\subsection{Constant Construction Mode Codes}

Constant construction instructions use opcode \code{0x08} and the U-type
instruction format.

\begin{center}
\begin{tabular}{lll}
\hline
\code{mode} value & Mnemonic & Meaning \\
\hline
\code{0x00} & \code{lui} & Load upper immediate, clear lower half \\
\code{0x01} & \code{lli} & Load lower immediate, clear upper half \\
\code{0x02} & \code{lhi} & Load upper immediate, preserve lower half \\
\code{0x03--0x3F} & reserved & Reserved \\
\hline
\end{tabular}
\end{center}

\noindent
For \code{lui}, the \code{imm16} field is placed in bits 31--16 of the
destination register and bits 15--0 are cleared.

\noindent
For \code{lli}, the \code{imm16} field is placed in bits 15--0 of the
destination register and bits 31--16 are cleared.

\noindent
For \code{lhi}, the \code{imm16} field is placed in bits 31--16 of the
destination register, without modifying bits 15--0.

\section{Memory Encoding Values}

Memory instructions use the M-type instruction format. Load instructions use
opcode \code{0x09}. Store instructions use opcode \code{0x0A}.

\subsection{Load Encoding}

Load instructions use opcode \code{0x09}. The \code{mode} field selects whether
sub-word loads are sign-extended or zero-extended; whether word loads are
integer or binary32; whether double loads are upper or lower.

\begin{center}
\begin{tabular}{llll}
\hline
Mnemonic & \code{size} & \code{mode} & Meaning \\
\hline
\code{lbu} & \code{00} & \code{0} & Load byte and zero-extend to 32 bits \\
\code{lb}  & \code{00} & \code{1} & Load byte and sign-extend to 32 bits \\
\code{lhu} & \code{01} & \code{0} & Load halfword and zero-extend to 32 bits \\
\code{lh}  & \code{01} & \code{1} & Load halfword and sign-extend to 32 bits \\
\code{lw}  & \code{10} & \code{0} & Load word \\
\code{ld}  & \code{10} & \code{1} & Load binary32 \\
\code{ldl}  & \code{11} & \code{0} & Load lower binary64 \\
\code{ldu}  & \code{11} & \code{1} & Load upper binary64 \\
\hline
\end{tabular}
\end{center}

\noindent
For byte and halfword loads, the loaded value is extended to 32 bits according
to the \code{mode} field.
\code{ld} extends upon loading from binary32 in memory to binary64 in register,
and loads into a floating-point register.
\code{ldl} and \code{ldu} load unchanged into their respective floating-point
register.

\subsection{Store Encoding}

Store instructions use opcode \code{0x0A}.

\begin{center}
\begin{tabular}{llll}
\hline
Mnemonic & \code{size} & \code{mode} & Meaning \\
\hline
\code{sb} & \code{00} & \code{0} & Store byte \\
\code{-} & \code{00} & \code{1} & Reserved \\
\code{sh} & \code{01} & \code{0} & Store halfword \\
\code{-} & \code{01} & \code{1} & Reserved \\
\code{sw} & \code{10} & \code{0} & Store word \\
\code{sd} & \code{10} & \code{1} & Store binary32 \\
\code{sdl} & \code{11} & \code{0} & Store lower binary64 \\
\code{sdu} & \code{11} & \code{1} & Store upper binary64 \\
\hline
\end{tabular}
\end{center}

\noindent
Executing an invalid store instruction raises an
\code{IllegalInstructionException}.

\subsection{Memory Access Size Field}

The \code{size} field has the same basic interpretation for load and store
instructions.

\begin{center}
\begin{tabular}{lll}
\hline
\code{size} value & Access size & Alignment requirement \\
\hline
\code{00} & Byte & Any address \\
\code{01} & Halfword & 2-byte aligned address \\
\code{10} & Word & 4-byte aligned address \\
\code{11} & Double-half & 4-byte aligned address \\
\hline
\end{tabular}
\end{center}

\noindent
Misaligned halfword, word and double-half memory accesses raise a
\code{MisalignedAccessException}. Byte accesses are always naturally aligned.

\section{Inter-File Transfer}

Inter-file transfer instructions convert values between the general-purpose and
floating-point register files. They use the R-type instruction format and major
opcode \code{0x10}.

The interpretation of the source and destination operands is determined by the
transfer function. Integer conversions interpret a general-purpose register as
a signed or unsigned 32-bit integer. Floating-point format conversions
interpret a general-purpose register as an IEEE-754 binary32 bit pattern or a
floating-point register as an IEEE-754 binary64 value.

\begin{center}
\begin{tabular}{llll}
\hline
\code{file} & \code{func} value & Mnemonic & Meaning \\
\hline
\code{0x0} & \code{0x000} & \code{disf} &
Test whether a binary64 value is representable as binary32 \\
\code{0x1} & \code{0x001} & \code{utod} &
Unsigned 32-bit integer to binary64 \\
\code{0x1} & \code{0x002} & \code{itod} &
Signed 32-bit integer to binary64 \\
\code{0x1} & \code{0x003} & \code{ftod} &
Binary32 to binary64 \\
\code{0x0} & \code{0x001} & \code{dtou} &
Binary64 to unsigned 32-bit integer \\
\code{0x0} & \code{0x002} & \code{dtoi} &
Binary64 to signed 32-bit integer \\
\code{0x0} & \code{0x003} & \code{dtof} &
Binary64 to binary32 \\
\code{-} & \code{0x004--0xFFF} & reserved &
Reserved \\
\hline
\end{tabular}
\end{center}

\noindent
For conversion instructions with a destination register, \code{rd} identifies
the destination register and \code{ra} identifies the source register. The
register file associated with each operand is determined by the transfer
function. Unless otherwise specified, \code{rb} is unused and must be encoded
as zero.

The \code{utod} and \code{itod} instructions perform numeric integer-to-
floating-point conversion. The \code{dtou} and \code{dtoi} instructions perform
numeric floating-point-to-integer conversion.

The \code{ftod} instruction interprets the 32 bits held in its general-purpose
source register as an IEEE-754 binary32 value and converts that value to
binary64. The \code{dtof} instruction converts an IEEE-754 binary64 value to
binary32 and writes the resulting 32-bit IEEE-754 representation to its
general-purpose destination register.

Floating-point conversion instructions update the applicable
\code{NZEIUX} flags according to the conversion result. The integer-only
\code{C} and \code{V} flags are cleared unless otherwise specified.

The \code{disf} instruction performs a representability test and does not
transfer a value between register files. Its precise flag behaviour is defined
with the instruction semantics.

\section{Branch Condition Codes}

Branch instructions use a 4-bit \code{cond} field to select the branch
condition. The interpretation of this field depends on the branch instruction
type.

BF-type instructions interpret \code{cond} as a direct test of one of the
architectural status flags. BC-type instructions interpret \code{cond} as a
comparison condition applied directly to their register operands.

The two instruction types therefore use independent condition-code
assignments.

\subsection{Flag-Based Branch Conditions}

BF-type instructions use opcode \code{0x0B}. The \code{cond} field selects one
of the eight architectural condition flags and whether the branch tests that
flag for being set or clear.

BF-type instructions do not interpret the flags as the result of a signed,
unsigned, integer, or floating-point comparison. They simply test the current
flag state.

\begin{center}
\begin{tabular}{llll}
\hline
\code{cond} value & Suffix & Branch is taken when & Meaning \\
\hline
\code{0x0} & \code{ns} & \flagN{} is set   & Negative set \\
\code{0x1} & \code{nc} & \flagN{} is clear & Negative clear \\
\code{0x2} & \code{zs} & \flagZ{} is set   & Zero set \\
\code{0x3} & \code{zc} & \flagZ{} is clear & Zero clear \\
\code{0x4} & \code{cs} & \flagC{} is set   & Carry set \\
\code{0x5} & \code{cc} & \flagC{} is clear & Carry clear \\
\code{0x6} & \code{vs} & \flagV{} is set   & Overflow set \\
\code{0x7} & \code{vc} & \flagV{} is clear & Overflow clear \\
\code{0x8} & \code{es} & \flagE{} is set   & Arithmetic error set \\
\code{0x9} & \code{ec} & \flagE{} is clear & Arithmetic error clear \\
\code{0xA} & \code{is} & \flagI{} is set   & Infinity set \\
\code{0xB} & \code{ic} & \flagI{} is clear & Infinity clear \\
\code{0xC} & \code{us} & \flagU{} is set   & Underflow set \\
\code{0xD} & \code{uc} & \flagU{} is clear & Underflow clear \\
\code{0xE} & \code{xs} & \flagX{} is set   & Inexact set \\
\code{0xF} & \code{xc} & \flagX{} is clear & Inexact clear \\
\hline
\end{tabular}
\end{center}

\noindent
The condition-code suffix is appended to the \code{bf} mnemonic. For example,
\code{bf.zs label} branches when the zero flag is set, while
\code{bf.ec label} branches when the arithmetic-error flag is clear.

\noindent
BF-type instructions test only the existing processor-status flags. They do
not perform a comparison and do not update the status register.

\subsection{Register-Comparison Branch Conditions}

BC-type instructions use opcode \code{0x0C} for the GPR, and \code{0x11} for the
FPR. The \code{cond} field selects a direct comparison between \code{ra} and
\code{rb}.

\begin{center}
\begin{tabular}{lll}
\hline
Suffix & Branch is taken when & Comparison type \\
\hline
\code{eq}  & \code{ra == rb} & Equality \\
\code{ne}  & \code{ra != rb} & Inequality \\
\code{lt}  & \code{ra < rb}  & Signed \\
\code{le}  & \code{ra <= rb} & Signed \\
\code{gt}  & \code{ra > rb}  & Signed \\
\code{ge}  & \code{ra >= rb} & Signed \\
\code{ltu} & \code{ra < rb}  & Unsigned \\
\code{leu} & \code{ra <= rb} & Unsigned \\
\code{gtu} & \code{ra > rb}  & Unsigned \\
\code{geu} & \code{ra >= rb} & Unsigned \\
\hline
\end{tabular}
\end{center}

\noindent
BC-type instructions do not read or update the condition flags. Condition-code
values \code{0xA}--\code{0xF} are reserved for BC-type instructions.

\section{Jump and Call Encoding Values}

\subsection{PC-Relative Jump and Call Opcodes}

PC-relative jump and call instructions use the J-type instruction format.

\begin{center}
\begin{tabular}{lll}
\hline
Opcode & Mnemonic & Meaning \\
\hline
\code{0x0D} & \code{jmp}  & PC-relative unconditional jump \\
\code{0x0E} & \code{call} & PC-relative call using \code{\$15} as the link
register \\
\hline
\end{tabular}
\end{center}

\noindent
Both instructions use the signed \code{offset26} field. The encoded offset is
measured in 32-bit instruction words and is shifted left by two bits to form the
byte offset. The target address is relative to the address of the instruction
following the jump or call.

\noindent
For \code{call}, the return address is written to the architectural link
register, \code{\$15}.

\subsection{Register-Indirect Jump and Call Function Codes}

Register-indirect jump and call instructions use opcode \code{0x0F} and the
JR-type instruction format.

\begin{center}
\begin{tabular}{lll}
\hline
\code{func} value & Mnemonic & Meaning \\
\hline
\code{0x00000} & \code{jr}   & Jump to address in \code{target} \\
\code{0x00001} & \code{jalr} & Jump to address in \code{target} and write link
register \\
\code{0x00002--0x3FFFF} & reserved & Reserved \\
\hline
\end{tabular}
\end{center}

\noindent
For \code{jr}, the \code{rd/link} field is reserved and must be encoded as
\code{0} by conforming assemblers.

\noindent
For \code{jalr}, the \code{rd/link} field selects the destination link register.
The standard ABI link register is \code{\$15}.

\noindent
Register-indirect targets must be 4-byte aligned. A misaligned register-indirect
target raises a \code{MisalignedFetchException}.

\section{System and Control Encoding Values}

System and control instructions use opcode \code{0x00} and the X-type
instruction format. The 18-bit \code{sysfunc} field is divided into a system
operation field and a payload field.

\begin{center}
\begin{tabular}{llll}
\hline
Bits within \code{sysfunc} & Instruction bits & Field & Meaning \\
\hline
17--12 & 25--20 & \code{sysop} & System operation selector \\
11--0  & 19--8  & \code{payload} & Operation-specific payload \\
\hline
\end{tabular}
\end{center}

\noindent
The interpretation of \code{payload} depends on the selected \code{sysop}. For
instructions that do not use a payload, the \code{payload} field is reserved
and must be encoded as zero by conforming assemblers.

\subsection{System Operation Codes}

\begin{center}
\begin{tabular}{lll}
\hline
\code{sysop} value & Mnemonic & Meaning \\
\hline
\code{0x00} & \code{nop}     & No operation \\
\code{0x01} & \code{halt}    & Halt processor execution \\
\code{0x02} & \code{trap}    & Raise software trap \\
\code{0x03} & \code{syscall} & Raise system-call exception \\
\code{0x04} & \code{iret}    & Return from interrupt or exception \\
\code{0x05} & \code{ei}      & Enable maskable interrupts \\
\code{0x06} & \code{di}      & Disable maskable interrupts \\
\code{0x07} & \code{rdpc}    & Read program counter \\
\code{0x08} & \code{mrs}     & Move from control register \\
\code{0x09} & \code{msr}     & Move to control register \\
\code{0x0A--0x3F} & reserved & Reserved \\
\hline
\end{tabular}
\end{center}

\subsection{System Payload Encoding}

\begin{center}
\begin{tabular}{llll}
\hline
Mnemonic & \code{sysop} & \code{payload} use & Payload requirement \\
\hline
\code{nop}     & \code{0x00} & Unused & Must be zero \\
\code{halt}    & \code{0x01} & Unused & Must be zero \\
\code{trap}    & \code{0x02} & \code{trapimm12} & Software-defined trap
immediate \\
\code{syscall} & \code{0x03} & Unused & Must be zero \\
\code{iret}    & \code{0x04} & Unused & Must be zero \\
\code{ei}      & \code{0x05} & Unused & Must be zero \\
\code{di}      & \code{0x06} & Unused & Must be zero \\
\code{rdpc}    & \code{0x07} & Unused & Must be zero \\
\code{mrs}     & \code{0x08} & \code{creg4} & Bits 11--4 must be zero \\
\code{msr}     & \code{0x09} & \code{creg4} & Bits 11--4 must be zero \\
\hline
\end{tabular}
\end{center}

\noindent
For \code{trap}, the 12-bit \code{payload} field is interpreted as an unsigned
software trap immediate, \code{trapimm12}. This data is stored in \code{\%edata}.

\noindent
For \code{mrs} and \code{msr}, the low 4 bits of \code{payload} encode the
selected control register. The upper 8 bits of \code{payload} are reserved and
must be encoded as zero by conforming assemblers.

\subsection{System Instruction Register Fields}

\begin{center}
\begin{tabular}{llll}
\hline
Mnemonic & \code{rd} & \code{rs} & Register-field use \\
\hline
\code{nop}     & \code{r0} & \code{r0} & Unused \\
\code{halt}    & \code{r0} & \code{r0} & Unused \\
\code{trap}    & \code{r0} & \code{r0} & Unused \\
\code{syscall} & \code{r0} & \code{r0} & Unused \\
\code{iret}    & \code{r0} & \code{r0} & Unused \\
\code{ei}      & \code{r0} & \code{r0} & Unused \\
\code{di}      & \code{r0} & \code{r0} & Unused \\
\code{rdpc}    & destination & \code{r0} & \code{rd} receives the program
counter value \\
\code{mrs}     & destination & \code{r0} & \code{rd} receives the selected
control-register value \\
\code{msr}     & \code{r0} & source & \code{rs} supplies the value written to
the selected control register \\
\hline
\end{tabular}
\end{center}

\noindent
For instructions that do not use \code{rd} or \code{rs}, unused register fields
must be encoded as \code{r0} by conforming assemblers.

\subsection{Control Register Numbers}

Control-register access instructions use the low 4 bits of the X-type
\code{payload} field to select a control register.

\begin{center}
\begin{tabular}{lll}
\textbf{creg4 value} & \textbf{Name} & \textbf{Meaning} \\
\code{0x0} & \code{\%pc}     & Program counter, read-only \\
\code{0x1} & \code{\%sr}     & Status register \\
\code{0x2} & \code{\%epc}    & Exception program counter \\
\code{0x3} & \code{\%esr}    & Exception status register \\
\code{0x4} & \code{\%ecause} & Exception cause \\
\code{0x5} & \code{\%edata}  & Exception-specific data \\
\code{0x6} & \code{\%evbase} & Exception vector base address \\
\code{0x7--0xF} & reserved & Reserved \\
\end{tabular}
\end{center}

\noindent
Executing \code{mrs} or \code{msr} with a reserved \code{creg4} value raises an
\code{IllegalInstructionException}.

\section{NOP Encoding}

VIV-32 defines the all-zero instruction word as \code{nop}.

\begin{center}
\begin{tabular}{lll}
\hline
Field & Value & Meaning \\
\hline
\code{opcode} & \code{0x00} & X-type system/control instruction \\
\code{sysop} & \code{0x00} & \code{nop} \\
\code{payload} & \code{0x000} & Unused \\
\code{rd} & \code{r0} & Unused \\
\code{rs} & \code{r0} & Unused \\
\hline
\end{tabular}
\end{center}

\noindent
The complete encoding of \code{nop} is therefore:

\begin{center}
\code{0x00000000}
\end{center}

\noindent
This guarantee allows zero-filled memory to decode as a valid no-operation
instruction. It also provides a safe default encoding for padding and alignment.

\section{Flag Update Policy}

Instruction encodings do not contain a general flag-update bit in VIV-32
v\version{}. Whether an instruction updates condition flags is defined by the
instruction semantics in Chapter~5.

Arithmetic, logical, comparison, shift, bit-test, multiply, divide and load
instructions update flags as specified by their instruction definitions.
Store, branch, jump, call, and most system/control instructions do not update
flags unless explicitly stated.

\section{Reserved Encodings}

Reserved opcode values, function codes, mode codes, condition codes, size
values, system operation codes, payload bits, control-register numbers, and
register-field combinations must not be emitted by conforming assemblers.

Unless otherwise specified, executing an instruction with a reserved encoding
raises an illegal-instruction exception.
